\begin{proof}
由 1.16 立见
$$LI_{ww'}(D(I(1,w'^{-1}\lambda)))=D(I(w'^{-1}w^{-1},w\lambda)).$$
另一方面，由 L.3.18 有
\begin{align*}
LI_{ww'}(D(I(1,w'^{-1}\lambda)))&=LI_w\bigl(LI_{w'}(D(I(1,w'^{-1}\lambda)))\bigr)\\
&=LI_w(D(I(w'^{-1},\lambda))).\qedhere
\end{align*}
\end{proof}

特别地，设 $w_0$ 是 $W$ 中的最长元，$w\in W$ 任意。则元素 $w'=w_0w$ 满足
$$ww'^{-1}=w_0^{-1}=w_0,\qquad\ell(w')=\ell(w_0)-\ell(w)$$
（见 [5], Ch.~VI, \S1, no.~6, Cor.~3. of Prop.~17）。由此
$$LI_{w'}(D(I(w,\lambda)))=D(I(w_0,w'\lambda)).$$

下一个结果对 1.14 的证明至关重要。

\begin{lemma}\label{lem:05.1.18}
设 $\lambda\in\h^*$。则
\begin{enumerate}[label=(\roman*)]
\item 当 $p>0$ 时 $H^p(X,I(w_0,\lambda))=0$；
\item $\Gamma(X,I(w_0,\lambda))=I(w_0\lambda)$。
\end{enumerate}
\end{lemma}

\begin{proof}
因 $C(w_0)$ 是 $X$ 的仿射开子簇，当 $p>0$ 时
$$H^p(X,I(w_0,\lambda))=H^p(C(w_0),\calO_{C(w_0)})=0.$$
这证明 (i)。

现证 (ii)。设 $\lambda\in\h^*$ 反支配。则由 1.15 有 $\Gamma(X,I(1,\lambda))=M(\lambda)$，且它是不可约 $\g$-模。取非零的 $v\in\Gamma(X,I(1,\lambda))$，它生成一个有限维 $\bb_0$-不变子空间 $U$。由 Engel 定理，存在 $v'\in U$，它张成一个 $\bb_0$-不变子空间。因此 $v'$ 是 $M(\lambda)$ 的权向量，对应某个权 $\mu\in\h^*$，且被 $\n_0$ 零化。于是存在从 $M(\mu+\rho)$ 到 $M(\lambda)$ 的自然态射。因 $M(\lambda)$ 不可约，有
$$L(\mu+\rho)=M(\lambda)=L(\lambda).$$
由 1.6 最终得到 $\lambda=\mu+\rho$。这样我们就证明了：$M(\lambda)$ 的每个 $\bb_0$-不变子空间 $U$ 都含最高权子空间 $M(\lambda)^{\lambda-\rho}$。

由 1.16 还可推出，对反支配的 $\lambda\in\h^*$ 有
$$\Gamma(X,I(w_0,w_0\lambda))=M(\lambda),$$
且它的每个 $\bb_0$-不变子空间都含最高权子空间。如前所述，对任意 $\mu\in\h^*$ 有
$$\Gamma(X,I(w_0,\mu))=\Gamma(C(w_0),\calO_{C(w_0)}).$$
因此 $C(w_0)$ 上的常函数 $1$ 是 $I(w_0,\mu)$ 的整体截面，它显然 $N_0$-不变。记 $x_{w_0}\in C(w_0)$ 为对应于含 $\h_0$ 且与 $\bb_0$ 相对的 Borel 子代数的点。$U^\circ|_{C(w_0)}$ 的截面
$$n\cdot x_{w_0}\longmapsto(\operatorname{Ad} n)\xi,\qquad \xi\in\h_0,$$
映为 $\D_{w_0\mu}|_{C(w_0)}$ 中的常截面 $(w_0\mu-\rho)(\xi)$，故它以 $(w_0\mu-\rho)(\xi)$ 乘截面 $1$。这又说明 $\xi$ 以 $(w_0\mu-\rho)(\xi)$ 乘此截面，即
$$1\in\Gamma(X,I(w_0,\mu))^{w_0\mu-\rho}.$$
特别地，若 $\lambda\in\h^*$ 反支配，则 $\Gamma(X,I(w_0,w_0\lambda))$ 的任一 $\bb_0$-不变子空间都含最高权子空间 $\Gamma(X,I(w_0,w_0\lambda))^{\lambda-\rho}$，它由 $C(w_0)$ 上的常函数组成。$I(w_0,\mu)$ 的几何平移是 $I(w_0,\mu)\otimes_{\calO_X}\calO(\nu)=I(w_0,\mu+\nu)$，又作为 $\calO_{C(w_0)}$-模有 $\calO(\nu)|_{C(w_0)}=\calO_{C(w_0)}$，故有 $\bb_0$-模的自然同构
$$\Gamma(X,I(w_0,\mu+\nu))\cong\Gamma(X,I(w_0,\mu))\otimes_\C\C_\nu.$$
于是可见，对任意 $\mu\in\h^*$：
\begin{enumerate}[label=(\alph*)]
\item 对任意权 $\omega\in\h^*$ 有
$$\dim\Gamma(X,I(w_0,\mu))^\omega=\dim M(w_0\mu)^\omega=\dim I(w_0\mu)^\omega;$$
\item $\Gamma(X,I(w_0,\mu))$ 的任一 $\bb_0$-不变子空间 $U$ 都含最高权子空间 $\Gamma(X,I(w_0,\mu))^{w_0\mu-\rho}$，它由 $C(w_0)$ 上的常函数组成。
\end{enumerate}
特别地，由 (b) 可知 $\Gamma(X,I(w_0,\mu))$ 的任一 $\g$-子模都含这些常函数，因此它只有一个不可约 $\g$-子模 $L$。再由 (a) 可知，$w_0\mu-\rho$ 是 $L$ 的最高权，即 $L=L(w_0\mu)$。最后由 1.12.(ii) 得到
$$\Gamma(X,I(w_0,\mu))=I(w_0\mu).$$
\end{proof}

现在可以证明 1.14。设 $\lambda\in\h^*$ 反支配，$w\in W$ 任意，且 $ww'^{-1}=w_0$。由上述讨论与 L.3.23 有
$$R\Gamma(D(I(w,\lambda)))=R\Gamma(LI_{w'}(D(I(w,\lambda))))=R\Gamma(D(I(w_0,w'\lambda))).$$
这推出
$$\Gamma(X,I(w,\lambda))=\Gamma(X,I(w_0,w'\lambda))=I(w_0w'\lambda)=I(w\lambda),$$
从而 1.14 得证。

现在确定不可约模的整体截面已相当直接。设 $W(\lambda)$ 是 $\lambda$ 的稳定化子。

\begin{theorem}\label{thm:05.1.19}
设 $\lambda\in\h^*$ 反支配。
\begin{enumerate}[label=(\roman*)]
\item 任一左 $W(\lambda)$-陪集都含最短元。
\item 设 $w\in W$。则下列断言等价：
\begin{enumerate}[label=(\alph*)]
\item $\Gamma(X,L(w,\lambda))\neq0$；
\item $\Gamma(X,L(w,\lambda))=L(w\lambda)$；
\item $w$ 是某个左 $W(\lambda)$-陪集中的最短元。
\end{enumerate}
\end{enumerate}
\end{theorem}

\begin{proof}
设 $w\in W$。首先，由 L.4.1 知 $\Gamma(X,L(w,\lambda))$ 或为不可约 $\g$-模，或为零。由 $\Gamma$ 的正合性与 1.14，若它非零，则它必是 $I(w\lambda)=\Gamma(X,I(w,\lambda))$ 的唯一不可约子模 $L(w\lambda)$。

设 $w$ 与 $w'$ 属于同一左 $W(\lambda)$-陪集，且 $\Gamma(X,L(w,\lambda))\neq0$、$\Gamma(X,L(w',\lambda))\neq0$。则 $w\lambda=w'\lambda$，从而 $L(w\lambda)=L(w'\lambda)$。由 L.4.2，这推出 $L(w,\lambda)\cong L(w',\lambda)$ 且 $w=w'$。因此，在一个左 $W(\lambda)$-陪集中，使 $\Gamma(X,L(w,\lambda))\neq0$ 的元素 $w$ 是唯一的。

设 $v\neq w$ 与 $w$ 同属一个左 $W(\lambda)$-陪集。则 $\Gamma(X,L(v,\lambda))=0$。由正合列
$$0\longrightarrow L(v,\lambda)\longrightarrow I(v,\lambda)\longrightarrow I(v,\lambda)/L(v,\lambda)\longrightarrow0$$
与 $\Gamma$ 的正合性可知
$$\Gamma(X,I(v,\lambda)/L(v,\lambda))=\Gamma(X,I(v,\lambda))=I(v\lambda)=I(w\lambda),$$
故后者以 $L(w\lambda)$ 为合成因子。这又说明 $L(w,\lambda)$ 是 $I(v,\lambda)/L(v,\lambda)$ 的合成因子。于是 $C(w)\subset\partial C(v)$ 且 $\ell(w)<\ell(v)$。因此 $\ell$ 在陪集 $wW(\lambda)$ 上于 $w$ 处取到最小值。
\end{proof}

特别地，若 $\lambda$ 还正则，则有下述结果。

\begin{corollary}\label{cor:05.1.20}
设 $\lambda\in\h^*$ 正则反支配。则对任意 $w\in W$ 有
$$\Gamma(X,L(w,\lambda))=L(w\lambda).$$
\end{corollary}

\begin{proof}
因 $\lambda$ 正则，$W(\lambda)=\{1\}$。
\end{proof}

1.14 的证明中的讨论有下述推论。

\begin{theorem}\label{thm:05.1.21}
设 $\lambda\in\h^*$，$\theta=W\cdot\lambda$。则 $M(\lambda)$ 在 $U_\theta$ 中的零化子是 $\{0\}$。
\end{theorem}

\begin{proof}
只需证明：对任意 $\lambda\in\h^*$，$I(\lambda)$ 的零化子都是平凡的。但由 1.18.(ii)，这等价于证明 $\Gamma(X,\D_\lambda)$ 中没有非平凡元素零化 $\Gamma(X,I(w_0,\lambda))=\Gamma(C(w_0),\calO_{C(w_0)})$，而这是显然的。
\end{proof}

最后我们讨论标准模 $I(w,\lambda)$ 与 Verma 模不可约的充要条件。先分析一个关键的特殊情形。

\begin{lemma}\label{lem:05.1.22}
设 $\lambda\in\h^*$，$\alpha\in\Pi$ 满足 $p=-\alpha^\vee(\lambda)\in\Z$。设 $w=w's_\alpha$ 且 $\ell(w)=\ell(w')+1$。则：
\begin{enumerate}[label=(\roman*)]
\item 有正合列
$$0\longrightarrow U^0(I(w',s_\alpha\lambda))\longrightarrow I(w,\lambda)\longrightarrow I(w',\lambda)\longrightarrow0;$$
\item $U^0(I(w',s_\alpha\lambda))\neq0$，且它是形如 $p_\alpha^+(V)$ 的模的平移。
\end{enumerate}
\end{lemma}

\begin{proof}
由 1.16 与 L.3.18 有
\begin{align*}
I(w,\lambda)&=I_{w^{-1}}(I(1,w\lambda))=I_{s_\alpha}\bigl(I_{w'^{-1}}(I(1,w\lambda))\bigr)\\
&=I_{s_\alpha}(I(w',s_\alpha\lambda)).
\end{align*}
于是由 L.5.3.(ii) 有正合列
\begin{center}
{\small
\begin{tikzcd}[column sep=small]
0 \arrow[r] & U^0(I(w',s_\alpha\lambda)) \arrow[r] & I(w,\lambda) \arrow[r] & I(w',s_\alpha\lambda)(-p_\alpha) \arrow[r] \arrow[d,"\cong"] & U^1(I(w',s_\alpha\lambda)) \arrow[r] & 0\\
& & & I(w',\lambda)
\end{tikzcd}
}
\end{center}
设 $p_\alpha:X\to X_\alpha$ 是旗簇 $X$ 到全体 $\alpha$ 型抛物子代数构成的簇的自然投影。则用 L.5 的记号，有交换图表
\begin{center}
\begin{tikzcd}
Y_\alpha \arrow[r,"q_2"] \arrow[d,"q_1"'] & X \arrow[d,"p_\alpha"]\\
X \arrow[r,"p_\alpha"'] & X_\alpha
\end{tikzcd}
\end{center}
由基变换，并利用 $p_\alpha\circ i_{w'}$ 是仿射簇 $C(w')$ 到 $X_\alpha$ 的浸入这一事实，得到 $U^1(I(w',s_\alpha\lambda))=0$；又 $U^0(I(w',s_\alpha\lambda))\neq0$，且它是形如 $p_\alpha^+(V)$ 的模的平移。这就推出两个断言。
\end{proof}

\begin{theorem}\label{thm:05.1.23}
设 $\lambda\in\h^*$，$w\in W$。则下列条件等价：
\begin{enumerate}[label=(\roman*)]
\item $\Sigma_w^+\cap\Sigma_\lambda=\emptyset$；
\item $I(w,\lambda)$ 是不可约 $\D_\lambda$-模。
\end{enumerate}
\end{theorem}

\begin{proof}
(i)$\Rightarrow$(ii) 若 $\Sigma_w^+\cap\Sigma_\lambda=\emptyset$，则由 L.3.22，交织函子 $I_w:\mathcal{M}_{qc}(\D_\lambda)\to\mathcal{M}_{qc}(\D_{w\lambda})$ 是范畴等价，$I_{w^{-1}}$ 是其逆。由 1.16 有
$$I_{w^{-1}}(I(1,w\lambda))=I(w,\lambda).$$
因 $I(1,w\lambda)$ 显然不可约，$I(w,\lambda)$ 也是不可约 $\D_\lambda$-模。

下面按 $\ell(w)$ 归纳，证明 $\Sigma_w^+\cap\Sigma_\lambda\neq\emptyset$ 时 $I(w,\lambda)$ 是可约 $\D_\lambda$-模。若 $\ell(w)=0$，则 $w=1$，断言显然。故可设断言对一切满足 $\ell(w')<k$ 的 $w'\in W$ 成立。设 $\ell(w)=k$。则存在 $\alpha\in\Pi$ 与 $w'\in W$，使 $w=w's_\alpha$ 且 $\ell(w')=k-1$。如上一证明，由 1.16 与 L.3.18 可得
$$I(w,\lambda)=I_{s_\alpha}(I(w',s_\alpha\lambda)).$$
又由 L.3.12.(ii)，
$$\Sigma_w^+\cap\Sigma_\lambda=s_\alpha\bigl(\Sigma_{w'}^+\cap\Sigma_{s_\alpha\lambda}\bigr)\cup\bigl(\{\alpha\}\cap\Sigma_\lambda\bigr).$$
若 $\alpha\notin\Sigma_\lambda$，则 $\operatorname{Card}(\Sigma_w^+\cap\Sigma_\lambda)=\operatorname{Card}(\Sigma_{w'}^+\cap\Sigma_{s_\alpha\lambda})$，故由归纳假设，$I(w',s_\alpha\lambda)$ 是可约 $\D_{s_\alpha\lambda}$-模。而由 L.3.22，此时 $I_{s_\alpha}:\mathcal{M}_{qc}(\D_{s_\alpha\lambda})\to\mathcal{M}_{qc}(\D_\lambda)$ 是范畴等价，故 $I(w,\lambda)$ 是可约 $\D_\lambda$-模。

若 $\alpha\in\Sigma_\lambda$，则由 1.22，$I(w,\lambda)$ 可约。
\end{proof}
\begin{annote}
底本此处把 1.22 引作 21，1.24 与 1.25 的证明中各有一处同样的引法。1.21 讲的是 $M(\lambda)$ 在 $U_\theta$ 中的零化子，既没有 (i) 款，也不含此处所需的 $\D_{w^{-1}\lambda}$-模正合列；给出该正合列的是 1.22。故三处均按 1.22 读出。[来源：底本 p.102，1.23 的证明；1.21 与 1.22 的陈述对勘]
\end{annote}

现在给出 Verma 模不可约的充要条件。

\begin{theorem}\label{thm:05.1.24}
设 $\lambda\in\h^*$。则下列条件等价：
\begin{enumerate}[label=(\roman*)]
\item $\lambda$ 反支配；
\item $M(\lambda)$ 不可约。
\end{enumerate}
\end{theorem}

\begin{proof}
(i)$\Rightarrow$(ii) 设 $\lambda$ 反支配。由 1.14 有 $I(\lambda)=\Gamma(X,I(1,\lambda))$。又 $I(1,\lambda)$ 显然不可约。由 L.4.1，$I(\lambda)$ 是不可约 $\g$-模，故 $M(\lambda)=\widetilde{I(\lambda)}$ 也不可约。

(ii)$\Rightarrow$(i) 取非反支配的 $\lambda$，并设 $M(\lambda)$ 不可约。设 $w\in W$ 是使 $w^{-1}\lambda$ 反支配的最短元。则由 1.14 有
$$\Gamma(X,I(w,w^{-1}\lambda))=I(\lambda)=\widetilde{M(\lambda)}=\widetilde{L(\lambda)}=L(\lambda)=M(\lambda).$$
取 $\alpha\in\Pi$，使 $w=w's_\alpha$ 且 $\ell(w)=\ell(w')+1$。则 $w'^{-1}\lambda$ 不是反支配的。我们断言 $p=-\alpha^\vee(w^{-1}\lambda)\in\N$。因 $w^{-1}\lambda$ 反支配，对任意 $\beta\in\Sigma^+$ 都有 $\beta^\vee(w^{-1}\lambda)\notin\N$。又 $s_\alpha$ 置换 $\Sigma^+-\{\alpha\}$ 中的根，故对任意 $\beta\in\Sigma^+-\{\alpha\}$ 有
$$(s_\alpha\beta)^\vee(w^{-1}\lambda)=\beta^\vee(w'^{-1}\lambda)\notin\N;$$
而 $w'^{-1}\lambda$ 不是反支配的，故 $\alpha^\vee(w'^{-1}\lambda)\in\N$。由 1.22 得到 $\D_{w^{-1}\lambda}$-模的正合列
$$0\longrightarrow U^0(I(w',w'^{-1}\lambda))\longrightarrow I(w,w^{-1}\lambda)\longrightarrow I(w',w^{-1}\lambda)\longrightarrow0.$$
因 $w^{-1}\lambda$ 反支配，由 C.3.2 得到正合列
$$0\to\Gamma(X,U^0(I(w',w'^{-1}\lambda)))\to\Gamma(X,I(w,w^{-1}\lambda))\to\Gamma(X,I(w',w^{-1}\lambda))\to0,$$
且 $\Gamma(X,I(w,w^{-1}\lambda))=I(\lambda)=\widetilde{M(\lambda)}$ 不可约。由 1.14 有
$$\Gamma(X,I(w',w^{-1}\lambda))=I(s_{w'\alpha}\lambda)\neq0.$$
这说明 $\Gamma(X,U^0(I(w',w'^{-1}\lambda)))=0$ 且 $I(\lambda)=I(s_{w'\alpha}\lambda)$。取对偶即知 $M(\lambda)=M(s_{w'\alpha}\lambda)$，而这只在 $\lambda=s_{w'\alpha}\lambda$ 时才可能。后者又说明 $s_\alpha$ 稳定 $w^{-1}\lambda$，与 $\alpha^\vee(w^{-1}\lambda)\in\N$ 矛盾。因此 $M(\lambda)$ 必可约。
\end{proof}

最后我们讨论单反射对应的交织函子对不可约模 $L(w,\lambda)$ 的作用。

\begin{proposition}\label{prop:05.1.25}
设 $\lambda\in\h^*$，$\alpha\in\Pi$。则：
\begin{enumerate}[label=(\roman*)]
\item 若 $\alpha^\vee(\lambda)\notin\Z$，则 $I_{s_\alpha}(L(w,\lambda))=L(ws_\alpha,s_\alpha\lambda)$。
\item 若 $\alpha^\vee(\lambda)\in\Z$，则有两种可能：
\begin{enumerate}[label=(\alph*)]
\item 若 $\ell(ws_\alpha)=\ell(w)+1$，则 $I_{s_\alpha}(L(w,\lambda))\neq0$，且 $L^{-1}I_{s_\alpha}(L(w,\lambda))=0$；
\item 若 $\ell(ws_\alpha)=\ell(w)-1$，则 $I_{s_\alpha}(L(w,\lambda))=0$，且 $L^{-1}I_{s_\alpha}(L(w,\lambda))=L(w,s_\alpha\lambda)$。
\end{enumerate}
\end{enumerate}
\end{proposition}

\begin{proof}
(i) 此时，由 L.3.22，$I_{s_\alpha}$ 是范畴等价，从而是正合函子。由 L.3.5 还可设 $\lambda$ 反支配且正则。于是 $s_\alpha\lambda$ 也反支配且正则。因此结论由 L.1.16、L.3.23 与 1.19 推出。

(ii) 设 $P_\alpha$ 是含 Borel 子群 $B$ 的 $\alpha$ 型抛物子群。则 $P_\alpha=B\cup Bs_\alpha B$。由此
$$p_\alpha^{-1}\bigl(p_\alpha(C(w))\bigr)=C(w)\cup C(ws_\alpha).$$
由于 $p_\alpha$ 是纤维同构于 $\mathbb{P}^1$ 的局部平凡投影，这个原像只能是二者之一：若 $\ell(ws_\alpha)=\ell(w)+1$，则 $C(w)\subset C(ws_\alpha)$，且 $p_\alpha^{-1}(p_\alpha(C(w)))=C(ws_\alpha)\neq C(w)$；若 $\ell(ws_\alpha)=\ell(w)-1$，则 $C(ws_\alpha)\subset C(w)$，且 $p_\alpha^{-1}(p_\alpha(C(w)))=C(w)$。因 $\Supp L(w,\lambda)=C(w)$，可见 $L(w,\lambda)$ 只有当 $\ell(ws_\alpha)=\ell(w)-1$ 时才可能是形如 $p_\alpha^+(V)$ 的模的平移。故由 L.5.6，$I_{s_\alpha}(L(w,\lambda))=0$ 蕴含 $\ell(ws_\alpha)=\ell(w)-1$。

下面证明逆命题。设 $\ell(ws_\alpha)=\ell(w)-1$。此时由 1.22 有正合列
$$0\longrightarrow U^0(I(ws_\alpha,s_\alpha\lambda))\longrightarrow I(w,\lambda)\longrightarrow I(ws_\alpha,\lambda)\longrightarrow0,$$
且 $U^0(I(ws_\alpha,s_\alpha\lambda))$ 是形如 $p_\alpha^+(V)$ 的模的非零平移。因此 $U^0(I(ws_\alpha,s_\alpha\lambda))$ 以 $L(w,\lambda)$ 为其唯一不可约子模，特别地
$$\Supp U^0(I(ws_\alpha,s_\alpha\lambda))=C(w).$$
又由 L.5.5，
$$L^{-1}I_{s_\alpha}\bigl(U^0(I(ws_\alpha,s_\alpha\lambda))\bigr)=U^0(I(ws_\alpha,\lambda)).$$
由 1.17，$L^{-1}I_{s_\alpha}(I(ws_\alpha,\lambda))$ 消没；于是对上述短正合列应用 $I_{s_\alpha}$ 的导出函子的长正合列，得到
$$L^{-1}I_{s_\alpha}(I(w,\lambda))=U^0(I(ws_\alpha,\lambda))\neq0.$$
又因 $L^{-1}I_{s_\alpha}$ 左正合，故 $L^{-1}I_{s_\alpha}(L(ws_\alpha,\lambda))=0$。现考虑短正合列
$$0\longrightarrow L(w,\lambda)\longrightarrow I(w,\lambda)\longrightarrow Q\longrightarrow0,$$
其中 $\Supp Q\subset\partial C(w)=\overline{C(w)}-C(w)$。由上述讨论，条件 $C(v)\subset\partial C(w)$ 与 $p_\alpha^{-1}(p_\alpha(C(v)))=C(w)$ 推出 $v=ws_\alpha$。因此，若 $L(v,\lambda)$ 是 $Q$ 的不可约成分且满足 $\Supp L^{-1}I_{s_\alpha}(L(v,\lambda))=C(w)$，则 $v=ws_\alpha$，而由上一注记这是不可能的。于是按 $Q$ 的长度归纳可知，$\Supp L^{-1}I_{s_\alpha}(Q)\neq C(w)$。因此，导出交织函子对应的长正合列中的那一段
\begin{center}
\begin{tikzcd}
0 \arrow[r] & L^{-1}I_{s_\alpha}(L(w,\lambda)) \arrow[r] & L^{-1}I_{s_\alpha}(I(w,\lambda)) \arrow[r] & L^{-1}I_{s_\alpha}(Q)\\
& & U^0(I(ws_\alpha,\lambda)) \arrow[u,"="']
\end{tikzcd}
\end{center}
说明第二个水平箭头非零，从而 $L^{-1}I_{s_\alpha}(L(w,\lambda))\neq0$。由 L.5.4，这推出 $I_{s_\alpha}(L(w,\lambda))=0$。于是 (ii) 是 L.5.4 的推论。
\end{proof}
